\section{Discussion}
\label{sec:discussion}

Benchmark F1 tracked KB yield for some design levers and not
others. Multi-corpus training and oncology-projected staging
improved out-of-domain benchmark generalisation, but those gains
did not translate uniformly to KB surfacing across encoder-schedule
cells. The disagreement concentrated in the training-schedule arm:
schedule explained 59.6\% of KB-argmax-accuracy variance but only
8.2\% of BioRED ex-NEG variance --- a roughly seven-fold gap.

Benchmark scores and KB audits measure overlapping but
non-identical constructs, and divergence under a specific design
lever is informative rather than disqualifying
\citep{lawlor2017triangulation,liugaunt2024asq}. BioRED measures
relation-type discrimination across a broad vocabulary on a topic
mix that is not oncology-specific; the CIViC audit measures
correctness on a narrow oncology subdomain with pre-identified
entity spans. Disagreement between the two scores does not
invalidate either; it shows that they answer different evaluation
questions and respond differently to specific design levers.

One plausible mechanism for the configuration-wave asymmetry is
sample-mix mismatch: CIViC targets concentrate in a narrow oncology
subdomain, whereas the BioRED test set spans a broader topic mix. Schedule
changes that add oncology-projected training shift model behaviour
in the subdomain the KB targets occupy densely but the benchmark
samples sparsely. This account is consistent with the seven-fold
schedule gap; we did not test it directly.

The inter-annotator audit (\S\ref{sec:results-iaa}) sharpens this
account. The heuristic projection over-attributes positive labels
to abstracts where one side of the CIViC-cited (drug, gene) pair is
not literally named, so KB argmax accuracy rewards positive
prediction even when the supporting evidence is not in the
abstract. All seven disagreements clustered on EGFR-pathway TKIs
and the HSP90 inhibitor tanespimycin. Oncology-projected training
plausibly increases positive predictions on these abstracts, so
part of the schedule lever's visibility on the KB metric may be
artefactual. This is a directional caveat; we did not test it
directly.

For clinical KB curation pipelines, the practical implication is to
treat benchmark F1 and a lightweight KB-anchor audit as two
parallel selection inputs. Configuration pairs within one
schema-wave BioRED SD of each other disagreed on KB argmax accuracy
by 0.16 points at the median, with half reversing their ranking, so
benchmark-alone selection is unreliable when the candidates differ
in training schedule. Training schedule therefore deserves more
attention than encoder scale for oncology curation pipelines. We
used 162 CIViC targets with pre-identified entity spans, requiring
no new annotation; the audit code and pre-registration are released
so that other curated oncology KBs (e.g., OncoKB, CancerMine) can
be substituted by replacing the source loader.

\paragraph{Limitations.}
The within-cell standard deviation of KB argmax accuracy (0.17)
is large relative to most paired contrasts we report; this widens
confidence intervals at our cell counts and renders the five
pre-registered coupling slopes inconclusive at the pre-registered
CI-width gate (\S\ref{sec:results-rq4}). The inter-annotator
agreement bound (\(\kappa=0.56\)) indicates moderate annotation
noise, and the seven disagreements concentrated on EGFR-pathway
TKIs would mechanically widen variance estimates in the schedule
arm. The second annotator was a large language model rather than an
independent human curator; \(\kappa=0.56\) should therefore be read
as a lower bound on agreement with an expert curator, and a
blinded human spot check on the same 30 targets is planned
follow-up. We evaluated against a single oncology KB;
generalisation to other clinical subdomains is untested. We
restricted to encoder-only models; an exploratory GPT-4o-mini KB baseline is summarised in Supplement~\ref{supp:llm-phase-d}.
There, GPT-4o-mini exceeds the always-DGR trivial baseline by only 0.04 on the 162-set, and reproduces the heuristic projector's label on all 7 IAA-disagreement targets; we read this as label-space artefact rather than evidence of LLM superiority on the underlying relation-discrimination task.
Other
follow-up work covers a wider LoRA sweep, multi-KB replication,
and a curator-adjudicated subset of the 162 CIViC targets to test
the directional-bias mechanism above.
